\documentclass[12pt]{article}
\usepackage[T1]{fontenc}
\usepackage[utf8]{inputenc}
\usepackage{lmodern}
\usepackage{textcomp}
\usepackage{enumitem}
\usepackage{geometry}
\geometry{margin=1in}
\begin{document}
\begin{center}
Answer Key - History - Undergraduate Level Assessment 16 \
\small (Answers are ordered exactly as in the question paper.)
\end{center}

\section*{Multiple Choice}
\begin{enumerate}[label=\arabic*.]
\item \textbf{Answer:} A
\\ \textit{Options:} A) focus on creating intricate geometric patterns.; B) construction of structures exclusively on hilltops.; C) employment of advanced hydraulic systems.; D) incorporation of wood, earth, and stone.; E) reliance on natural materials such as leaves and vines.
\\ \textit{Note:} Inca architecture is characterized by its remarkable attention to detail in creating intricate geometric patterns, which reflects their advanced engineering skills and cultural emphasis on symmetry and aesthetics in building design.
\item \textbf{Answer:} A
\\ \textit{Options:} A) Aurignacian; Solutrean; B) Gravettian; Solutrean; C) Gravettian; Aurignacian; D) Gravettian; Magdelanian; E) Magdelanian; Gravettian
\\ \textit{Note:} The Aurignacian tradition is known for its sophisticated stone tools, particularly finely made, leaf-shaped blades, while the subsequent Solutrean tradition is characterized by its advanced craftsmanship in bone and antler artifacts.
\item \textbf{Answer:} A
\\ \textit{Options:} A) meat, fish, and some plants.; B) primarily C$_{3}$ plants that included barley and oats.; C) large quantities of insects and small mammals.; D) a balanced diet of fruits, vegetables, meat, and fish.; E) primarily C$_{4}$ plants that included corn and potatoes.
\\ \textit{Note:} The analysis of the teeth of the hominid species from Malapa Cave reveals wear patterns and isotopic signatures consistent with a diet rich in meat and fish, supplemented by some plant materials, indicating a diverse and opportunistic feeding strategy.
\item \textbf{Answer:} A
\\ \textit{Options:} A) the foramen magnum; B) the skull; C) the feet; D) the femur; E) the spine
\\ \textit{Note:} The foramen magnum's position at the base of the skull indicates whether an organism's head is balanced above the spine for bipedal locomotion or further back for quadrupedal movement, making it a critical determinant in classifying primate locomotion.
\item \textbf{Answer:} C
\\ \textit{Options:} A) 100; B) 500; C) 1,000; D) 5,000
\\ \textit{Note:} The correct answer is 1,000 because archaeological evidence from the Minoan palace complex at Knossos suggests it was an extensive and complex structure with approximately 1,300 rooms, indicating its significance as a center of administration, culture, and economy in ancient Crete.
\end{enumerate}

\section*{True / False}
\begin{enumerate}[label=\arabic*.]
\setcounter{enumi}{5}
\item \textbf{Answer:} True
\\ \textit{Options:} A) True; B) False
\\ \textit{Note:} According to the passage: growth of private industry
\item \textbf{Answer:} False
\\ \textit{Options:} A) True; B) False
\\ \textit{Note:} The correct answer is: Causing a Commotion
\item \textbf{Answer:} False
\\ \textit{Options:} A) True; B) False
\\ \textit{Note:} The correct answer is: Al-Qaeda
\item \textbf{Answer:} True
\\ \textit{Options:} A) True; B) False
\\ \textit{Note:} According to the passage: spoken mainly in Thessaly, Boeotia and Lesbos
\item \textbf{Answer:} True
\\ \textit{Options:} A) True; B) False
\\ \textit{Note:} According to the passage: The East Slavs colonised Siberia and Central Asia.
\end{enumerate}

\section*{Long Form Response}
\begin{enumerate}[label=\arabic*.]
\setcounter{enumi}{10}
\item \textbf{Answer:} 1676
\\ \textit{Note:} Expected answer: 1676
\item \textbf{Answer:} the Phagmodrupa Dynasty
\\ \textit{Note:} Expected answer: the Phagmodrupa Dynasty
\end{enumerate}

\vfill
\noindent\textit{End of Answer Key}
\end{document}